
\svnid{$Id: $}
\svnid{$URL: $}

\section{FOR REVIEW:Precipitation on 1D}
\newrefsegment

%---------------------------------------------------------------------------------
\section*{Quality Assurance}
\begin{tabular}{@{}p{27mm}@{}p{0mm}p{\textwidth-52mm-48pt}p{15mm}p{10mm}}
\textbf{Date} & \textbf{:} &   \\ 
\textbf{Author} & \textbf{:} & Eskedar Gebremedhin  & \textbf{Initials:} & E.G. \\
\textbf{Review} & \textbf{:} &   & \textbf{Initials:} & \ldots 
\end{tabular}

%---------------------------------------------------------------------------------
\section*{Version information}
{{
\begin{tabular}{@{}p{27mm}@{}p{0mm}p{\textwidth-27mm-24pt}}
\textbf{Executable}   & \textbf{:} & Deltares, D-Flow FM Version 1.2.73.65109M, October 09 2019, 25:48:09  \\
\textbf{Location input}     & \textbf{:} & \url{https://repos.deltares.nl/repos/DSCTestbench/trunk/cases/e02_dflowfm/f102_lateral-flows/c02_1d-precipitation, https://repos.deltares.nl/repos/DSCTestbench/trunk/cases/e02_dflowfm/f102_lateral-flows/c02a_1d-precipitation,
https://repos.deltares.nl/repos/DSCTestbench/trunk/cases/e02_dflowfm/f102_lateral-flows/c02b_1d-precipitation, https://repos.deltares.nl/repos/DSCTestbench/trunk/cases/e02_dflowfm/f102_lateral-flows/c02c_1d-precipitation, https://repos.deltares.nl/repos/DSCTestbench/trunk/cases/e02_dflowfm/f102_lateral-flows/c02d_1d-precipitation,
https://repos.deltares.nl/repos/DSCTestbench/trunk/cases/e02_dflowfm/f102_lateral-flows/c02e_1d-precipitation} \\
\textbf{Revision input} & \textbf{:} & \\
\textbf{Location output}     & \textbf{:} & \url{https://repos.deltares.nl/repos/DSCTestbench/trunk/references/win64/e02_dflowfm/f102_lateral-flows/c02_1d-precipitation/DFM_OUTPUT_model1, https://repos.deltares.nl/repos/DSCTestbench/trunk/references/win64/e02_dflowfm/f102_lateral-flows/c02a_1d-precipitation/DFM_OUTPUT_model2,
https://repos.deltares.nl/repos/DSCTestbench/trunk/references/win64/e02_dflowfm/f102_lateral-flows/c02b_1d-precipitation/DFM_OUTPUT_model3,https://repos.deltares.nl/repos/DSCTestbench/trunk/references/win64/e02_dflowfm/f102_lateral-flows/c02c_1d-precipitation/DFM_OUTPUT_model4, 
https://repos.deltares.nl/repos/DSCTestbench/trunk/references/win64/e02_dflowfm/f102_lateral-flows/c02d_1d-precipitation/DFM_OUTPUT_model5, https://repos.deltares.nl/repos/DSCTestbench/trunk/references/win64/e02_dflowfm/f102_lateral-flows/c02e_1d-precipitation/DFM_OUTPUT_model6} \\
\textbf{Revision output} & \textbf{:} & 
\end{tabular}
}}

%---------------------------------------------------------------------------------
\section*{Purpose}
\textbf{FOR REVIEW} 
The purpose of this validation test is to verify the proper implementation of precipitation on 1D \dflowfm. 

%---------------------------------------------------------------------------------
\section*{Linked claims}
\dflowfm can use various methods to apply laterals on 1D.
%---------------------------------------------------------------------------------
\section*{Approach}
This test case includes six separate models, five of the models are with a straight river 10 km, various cross-section per model and the last model is a straight channel that split into two branches (each 5 km) at chainage=5000 m.
1 mm/hr constant precipitation is imposed to fill up the channels' storage area. The precipiation volume in \dflowfm is applied on the actual channel length and the storage volume and water level is distributed with an additional grid point (half a grid at both sides of the channel).

The analytical solution for the water level is derived from a calculation of the relationship of the storage (taking into account varying profile shape) in the systems and the total precipitation.

%---------------------------------------------------------------------------------
\section*{Conclusion} 
Accurate implementation of precipiation on cross-section profiles (rectangle, ZW and YZ) where it relate to storage volume is confirmed. 
The modelled water level and total volume extracted from all locations in the models are identical to the analytical solution which is calculated considering the \dflowfm additional grid point.

%---------------------------------------------------------------------------------
\section*{Model description}
The test comprises of 6 separate models (see \autoref{fig:e02_f102_c02_Test model lay-out}) with the following characteristics.

\begin{itemize}
	\item All models consist of a horizontal bed level at 0 m AD.
	\item T1-T5 are 10 km and T6 is with three branches each 5 km in length with an equidistant grid spacing of 500 m. 
	\item All models have a constant hydraulic roughness of Ch\'ezy = 45 m\textsuperscript{1/2}/s.    
	\item In all models, 1 mm/hr constant precipitation for 10 days period, 5 m initial water depth with closed boundaries is applied. 
	\item In T2, T3 and T4 left cross-section is used at chainage = 0.00 m and right cross-section is used at chainage = 10000.00 m 
	\item T1 and T6 is a model with a uniform rectangular cross-sectional profile with 100 m width and 10 m level.
	\item T2 is a model with varying rectangular cross-sectional profile - left: 100 m wide, right: 500 m wide and 10 m level.
	\item T3 is a model with varying ZW cross-sectional profile - left: shown in \autoref{fig:e02_f102_c02_zw_left}, right: shown in \autoref{fig:e02_f102_c02_zw_right}.
	\item T4 is a model with varying YZ cross-sectional profile - left: shown in \autoref{fig:e02_f102_c02_yz_left}, right: shown in \autoref{fig:e02_f102_c02_yz_right}.
	\item T5 is a model with a uniform YZ cross-sectional profile : shown in \autoref{fig:e02_f102_c02_yz}.
\end{itemize}

\begin{figure}[H]
    \centering
    \includegraphics*[width=0.95\textwidth]{figures/e02_f102_c02_layout.jpg}
    \caption{Layout of the test models, equidistant grid spacing at connection nodes}
    \label{fig:e02_f102_c02_Test model lay-out}
\end{figure} 

\begin{figure}[H]
    \centering
    \includegraphics*[width=0.95\textwidth]{figures/e02_f102_c02_T3_left_cross-section.jpg}
    \caption{T3 left side ZW-profile}
    \label{fig:e02_f102_c02_zw_left}
\end{figure}
 
\begin{figure}[H]
    \centering
    \includegraphics*[width=0.95\textwidth]{figures/e02_f102_c02_T3_right_cross-section.jpg}
    \caption{T3 right side ZW-profile}
    \label{fig:e02_f102_c02_zw_right}
\end{figure}

\begin{figure}[H]
    \centering
    \includegraphics*[width=0.95\textwidth]{figures/e02_f102_c02_T4_left_cross-section.jpg}
    \caption{T4 left side YZ-profile}
    \label{fig:e02_f102_c02_yz_left}
\end{figure}
 
\begin{figure}[H]
    \centering
    \includegraphics*[width=0.95\textwidth]{figures/e02_f102_c02_T4_right_cross-section.jpg}
    \caption{T4 right side YZ-profile}
    \label{fig:e02_f102_c02_yz_right}
\end{figure}

\begin{figure}[H]
    \centering
    \includegraphics*[width=0.95\textwidth]{figures/e02_f102_c02_T5_cross-section.jpg}
    \caption{T5 YZ-profile}
    \label{fig:e02_f102_c02_yz}
\end{figure}

%---------------------------------------------------------------------------------
\section*{Results}
Results of all models (T1 to T6) are presented in \autoref{tab:e02_f102_c02_result}. 

\begin{longtable}{|p{10mm}|p{20.5mm}|p{18.5mm}|p{11.5mm}|p{20.5mm}|p{18.5mm}|p{11.5mm}|p{18.5mm}|}
\caption [Precipitation volume, storage volume and water level in models T1-T6.] {Precipitation volume, storage volume and water level in models T1-T6.} 
\label{tab:e02_f102_c02_result} \\ 
\hline
\STRUT
{All}       & \multicolumn{3}{c|}{Modelled results}  &  \multicolumn{3}{c|}{Analytical results} &  {Water}  \\
\cline{2-4}
\cline{5-7}  
\STRUT
{Models}    & \multicolumn{2}{c|}{Total volume}   & {Water}   & \multicolumn{2}{c|}{Total volume}  & {Water}  & {level}  \\
\cline{2-3}
\cline{5-6}
\STRUT
			       & {precipitation \unitbrackets{m3}} & {storage \unitbrackets{m3}} & {level \unitbrackets{m AD}} & {precipitation \unitbrackets{m3}} & {storage \unitbrackets{m3}} & {level \unitbrackets{m AD}} & {Differences \unitbrackets{m AD}}  \\
[1ex] \hline \endhead
\hline
\endfoot
%
\STRUT
T1 		& 240000  & 5490000   & 5.22857 & 240000  & 5490000   & 5.22857 & 0.00 \\
\STRUT                                                                  
T2 		& 720000  & 16470000  & 5.22857 & 720000  & 16470000  & 5.22857 & 0.00 \\
\STRUT                                                                  
T3 		& 720000  & 5838750   & 5.40234 & 720000  & 5838750   & 5.40234 & 0.00 \\
\STRUT                                                                  
T4 		& 720000  & 4657500   & 5.43796 & 720000  & 4657500   & 5.43796 & 0.00 \\
\STRUT                                                                  
T5 		& 1200000 & 3824738   & 6.14295 & 1200000 & 3824738   & 6.14295 & 0.00 \\
\STRUT                                                                  
T6 		& 360000  & 8235000   & 5.22857 & 360000  & 8235000   & 5.22857 & 0.00 \\
[1ex] \hline
\end{longtable}

%---------------------------------------------------------------------------------
\section*{Analysis of results}
The channels are filled through a constant precipitation over 10 days. 
The total volume \unitbrackets{m3} of the models after 10 day is 40.56 milion m3 and the  water level have risen according to the shape of the profile.

A comparison of precipitation, water level and total storage volume in all models is in line with the analytical solution.

It should be noted that \dflowfm simulations adds an additional half grid spacing at the start and end of branch of the models.
As a result, the precipiation that is falling on the total length of the channel is distributed over the additional grid point. 

Moreover, in the case of T6 model, the three channel's full length is used in \dflowfm to compute the total volume which is double counting the storage area on the junction of the branches. 

These points should be taken into consideration when calculating the storage volume and  water level of the channels.

%---------------------------------------------------------------------------------
\printrefsegment
